\subsection{基本域上的积分}

设 X 是局部紧空间，H 是在 X 上连续且 逆紧地从右作用的离散群。令 \(\pi\) 为从 X 到 X/H 的典范映射。对于每个 \(x \in X\)，记 H 中 x 的稳定子为 \(H_x\)；这是 H 的有限子群（GT, III, §4, No.~2, Prop. 4）；其阶记为 \(n(x)\)。对于每个 \(s \in H\)，有 \(H_{xs} = s^{-1}H_x s\)，因而 \(n(xs) = n(x)\)。存在 x 的开邻域 U，使得 \(U \cap Us = \varnothing\) 对 \(s \notin H_x\) 成立（同上, No.~4, Prop. 8 的证明）；对于 \(y \in U\)，有 \(H_y \subset H_x\)；因此 X 上的函数 n 是上半连续的。当 X 在无穷远处可数时，H 可数；事实上，设 \((\mathrm{K}_1, \mathrm{K}_2, \ldots)\) 是由紧子集序列覆盖 X 的覆盖，并取 \(x_{0} \in X\)；\(s \in H\) 满足 \(x_{0}s \in K_{i}\) 的这些元素所成的集合是有限的（同上, No.~5, Th. 1），从而得到所述结论。

定义 2. --- 设 \(F \subset X\)。如果 \(\pi\) 在 F 上的限制是从 F 到 X/H 的双射（换言之，F 是 H 所定义的等价关系的代表元系），则称 F 为（关于 H 的）基本域。

引理 8. --- 设 F 是基本域。对于每个 \(x \in \mathbf{X}\)，

\[
\sum_ {s \in \mathrm{H}} \varphi_ {\mathrm{F} s} (x) = n (x).\tag{24}
\]

由于有 \(\varphi_{\mathrm{Fs}}(xt) = \varphi_{\mathrm{Fst}^{-1}}(x)\)，对所有 \(s\) 和 \(t\) 属于 \(\mathbf{H}\) 成立，当 \(x\) 替换为 \(xt\) 时，(24) 的两边仍保持不变。因此可设 \(x \in \mathbf{F}\)。此时有以下等价关系

\[
\varphi_ {\mathrm{Fs}} (x) = 1 \Leftrightarrow x \in \mathrm{Fs} \Leftrightarrow x s ^ {- 1} \in \mathrm{F} \Leftrightarrow x s ^ {- 1} = x \Leftrightarrow s \in \mathrm{H} _ {x},
\]

从而得到 (24)。

命题 15. --- 假设 X 在无穷远处可数。设 \(\mu\) 是 X 上的 \(\geqslant 0\) 测度。设 F 是基本域，并且 Fs 都是 \(\mu\)-可测的，对每个 \(s \in H\) 成立。设 f 是一个 \(\mu\)-可积函数，定义在 X 上，取值于 Banach 空间或 \(\overline{\mathbf{R}}\)。则族

\[
\int_ {\mathrm{Fs}} n (x) ^ {- 1} \mathbf {f} (x) d \mu (x) \quad (s \in \mathbf {H})
\]

是可求和的，并且

\[
\int_ {\mathrm{X}} \mathbf {f} (x) d \mu (x) = \sum_ {s \in \mathrm{H}} \int_ {\mathrm{F} s} n (x) ^ {- 1} \mathbf {f} (x) d \mu (x).
\]

若 \(\mathbf{A}\) 是 \(\mathbf{H}\) 的有限子集，则

\[
\left| \sum_ {s \in \mathrm{A}} n ^ {- 1} \mathbf {f} \varphi_ {\mathrm{Fs}} \right| \leqslant n ^ {- 1} | \mathbf {f} | \sum_ {s \in \mathrm{A}} \varphi_ {\mathrm{Fs}} \leqslant | \mathbf {f} |
\]

由引理 8 得出。引理 8 还证明，\(\sum_{s\in\mathbf{A}}n^{-1}\mathbf{f}\varphi_{\mathrm{F}s}\) 逐点收敛到 \(\mathbf{f}\)，其中指标集是由 \(\mathbf{H}\) 的有限子集构成的递增有向集。于是 Prop. 15 由 Ch. IV, §4, No.~3, Th. 2 得出。

定理 4. --- 设 X 是在无穷远处可数的局部紧空间，H 是在 X 上连续且 逆紧地从右作用的离散群，\(\pi\) 是从 X 到 X/H 的典范映射，\(\mu\) 是 X 上关于 H 不变的正测度，\(\beta\) 是 H 的归一化 Haar 测度，且 \(\lambda = \mu/\beta\)。设 F 是 \(\mu\)-可测的基本域。

\begin{enumerate}
\def\labelenumi{\alph{enumi})}
\tightlist
\item
  偶 \((\pi, n^{-1}\varphi_{\mathrm{F}})\) 对 \(\mu\) 是适合的，并且
\end{enumerate}

\[
\int_ {\mathrm{X}} n (x) ^ {- 1} \varphi_ {\mathrm{F}} (x) \varepsilon_ {\pi (x)} d \mu (x) = \lambda .
\]

\begin{enumerate}
\def\labelenumi{\alph{enumi})}
\setcounter{enumi}{1}
\item
  映射 \(\pi\) 关于 \(n^{-1}\varphi_{\mathrm{F}}\cdot \mu\) 是 逆紧的，并且 \(\pi (n^{-1}\varphi_{\mathrm{F}}\cdot \mu) = \lambda\)。
\item
  设 k 是 X/H 上的函数。k 关于 \(\lambda\) 可测（分别为 \(\lambda\)-可积）的充要条件是 \(n^{-1}\varphi_{\mathrm{F}}(k\circ\pi)\) 关于 \(\mu\) 可测（分别为 \(\mu\)-可积）；并且，若 k 关于 \(\lambda\) 可积，则
\end{enumerate}

\[
\int_ {\mathrm{X/H}} k d \lambda = \int_ {\mathrm{F}} n ^ {- 1} (k \circ \pi) d \mu .
\]

我们有 \(\mu = \lambda^{\sharp}\)。设 \(f \in \mathcal{H}_{+}(\mathrm{X / H})\)。则 \(n^{-1}\varphi_{\mathrm{F}}(f \circ \pi)\) 关于 \(\mu\) 可测且 \(\geqslant 0\)，由 No.~3 的 Prop. 5 b) 有

\[
\int_ {\mathrm{X}} ^ {*} n (x) ^ {- 1} \varphi_ {\mathrm{F}} (x) f (\pi (x)) d \mu (x) = \int_ {\mathrm{X/H}} ^ {*} f (\dot {x}) d \lambda (\dot {x}) \int_ {\mathrm{H}} ^ {*} n (x \xi) ^ {- 1} \varphi_ {\mathrm{F}} (x \xi) d \beta (\xi)
\]

并且由引理 8，\(\int_{\mathrm{H}}^{*}n(x\xi)^{-1}\varphi_{\mathrm{F}}(x\xi)d\beta (\xi) = n(x)^{-1}\sum_{\xi \in \mathrm{H}}\varphi_{\mathrm{F}}(x\xi) = 1\)。因此 \(n^{-1}\varphi_{\mathrm{F}}\cdot (f\circ \pi)\) 关于 \(\mu\) 可积，并且

\[
\int_ {\mathrm{X}} n (x) ^ {- 1} \varphi_ {\mathrm{F}} (x) f (\pi (x)) d \mu (x) = \int_ {\mathrm{X/H}} f (\dot {x}) d \lambda (\dot {x}).
\]

这证明了 a)。b) 的断言类似可证。c) 的断言可由 a) 以及 Ch. V, §4, Prop. 3 和 Th. 2 得出。

推论。--- 保持 Th. 4 的假设和记号。设 F' 是第二个 μ-可测基本域。设 u 是 X 上取值于 Banach 空间或 \(\overline{R}\) 的函数，并且关于 H 不变。假设 u 在 F 上 μ-可积。则 u 在 F' 上 μ-可积，并且

\[
\int_ {\mathrm{F}} \mathbf {u} (x) d \mu (x) = \int_ {\mathrm{F} ^ {\prime}} \mathbf {u} (x) d \mu (x).
\]

由于 \(\mathbf{u}\) 和 \(n\) 关于 \(\mathbf{H}\) 不变，存在函数 \(\mathbf{v}\)，定义在 \(\mathrm{X / H}\) 上，使得 \(\mathbf{v} \circ \pi\) 与 \(n\mathbf{u}\) 在 \(\mathbf{F}\) 和 \(\mathbf{F}'\) 上一致。于是 \(n^{-1}\varphi_{\mathrm{F}}(\mathbf{v}\circ \pi) = \varphi_{\mathrm{F}}\mathbf{u}, n^{-1}\varphi_{\mathrm{F}^{\prime}}(\mathbf{v}\circ \pi) = \varphi_{\mathrm{F}^{\prime}}\mathbf{u}.\) 根据假设，\(n^{-1}\varphi_{\mathrm{F}}(\mathbf{v}\circ \pi)\) 关于 \(\mu\) 可积。由 Th. 4，\(\mathbf{v}\) 关于 \(\lambda\) 可积，\(\varphi_{\mathrm{F}^{\prime}}\mu\) 关于 \(\mu\) 可积，并且

\[
\int_ {\mathrm{F}} \mathbf {u} d \mu = \int_ {\mathrm{X/H}} \mathbf {v} d \lambda = \int_ {\mathrm{F} ^ {\prime}} \mathbf {u} d \mu .
\]

关于 \(\mu\)-可测基本域的存在性，见 Exer. 12。

